\section{Layer I: Biological Grounding}

\subsection{Why Biology Belongs In A Chapter On Formation}

The neighbouring consciousness chapter established one strong empirical floor:
biological systems exhibit organized state signatures, stability structure, and
nontrivial transition dynamics. That result does not yet prove ``framework
formation'' in the everyday sense, but it changes what is intellectually
available. The book is now allowed to ask not only how states persist, but how
repeated coupling and conditioning build regimes that later alter transition
cost.

\begin{proposition}[History-sensitive inertia]
The inertia of a current state depends not only on the present stimulus, but
also on the accumulated reinforcement history, expected reward landscape, and
prior successful re-entry routes associated with that state.
\end{proposition}

\begin{discussion}
This proposition is not meant as a direct theorem extracted from one paper. It
is a synthesis claim justified at \textbf{medium} strength by the chapter's
sources. The empirical consciousness papers support taking structured states
seriously. The causal-emergence and metastability sources support treating
history and system organization as explanatory variables rather than merely
decorative background.
\end{discussion}

\subsection{Conditioning and Integrative Organization}

The 2025 Communications Biology paper on associative conditioning in gene
regulatory network models is especially useful because it gives a tractable
example of repeated local interaction producing stronger higher-level
organization. This source should be read carefully. It does \emph{not} show
that human study routines are literally gene regulatory networks. What it buys
the manuscript is more limited and more valuable:
\begin{enumerate}[nosep]
  \item repeated local couplings can alter global organization,
  \item higher-level organization can acquire explanatory value,
  \item conditioning history can change what later transitions look like.
\end{enumerate}

For framework formation, this means a regime is not merely the sum of discrete
past acts. Once cues, rewards, expectations, and environmental reinforcements
become mutually stabilizing, the resulting pattern becomes a real operating
layer. In business terms, the system begins to ``carry policy'' in its own
structure.

\subsection{Metastability Rather Than Simple Rigidity}

The 2024 Nature Reviews Neuroscience article on metastability is critical
because it blocks a false binary. Frameworks are often described as either
rigid habits or flexible choices. Metastability offers a better middle term:
organized patterns can be durable enough to persist and yet reconfigurable
enough to release under changed conditions.

\begin{definition}[Metastable framework]
A \textbf{metastable framework} is a regime that is locally persistent under
ordinary perturbations, but not globally fixed against larger contextual,
reward, or control-parameter changes.
\end{definition}

\begin{discussion}
This definition matters because many productive frameworks must be metastable,
not absolute. A study routine that can never be interrupted is dysfunctional.
A routine that collapses under every disturbance is useless. The chapter
therefore treats metastability as the right biological and systems-level bridge
for explaining why good frameworks feel stable without becoming cages.
\end{discussion}
